\documentclass{chsmc}
\chsmcsetup{
  year = 2015,
  part = 1,
  type = A,
  show-answers = false,
}
\begin{document}

\section{填空题：本大题共 8 小题，每小题 8 分，共 64 分.}

\begin{enumerate}
  \item 设 $a$, $b$ 为不相等的实数，若二次函数 $f(x) = x^2+ax+b$ 满足
    $f(a) = f(b)$，则 $f(2)$ 的值是 \blankline
    \answer{$4$}
    \begin{solution}
      由已知条件及二次函数图像的轴对称性，可得 $\frac{a+b}{2} = -\frac{a}{2}$，即 $2a+b = 0$，所以 $f(2) = 4+2a+b = 4$.
    \end{solution}
  \item 若实数 $\theta$ 满足 $\cos\theta = \tan\theta$,
    则 $\frac{1}{\sin\theta} + \cos^4\theta$ 的值为 \blankline
    \answer{$2$}
    \begin{solution}
      由条件知，$\cos^2\alpha$
    \end{solution}
  \item 已知复数数列 $\{z_n\}$ 满足 $z_1=1$, $z_{n+1} = \bar{z_n}+1+n\ii$ ($n=1,2,\ldots$)，其中 $\ii$ 为虚数单位，$\bar{z_n}$ 表示 $z_n$ 的共轭复数，则 $z_{2015}$ 的值是 \blankline
    \answer{$2015 + 2017\ii$}
  \item 在矩形 $ABCD$ 中，$AB=2$，$AD=1$，边 $DC$ 上（包括点 $D$, $C$）的动点 $P$ 与 $CB$ 延长线上（包括点 $B$）的动点 $Q$ 满足 $|\overrightarrow{DP}| = |\overrightarrow{BQ}|$，则向量 $\overrightarrow{PA}$ 与向量 $\overrightarrow{PQ}$ 的数量积 $\overrightarrow{PA}\cdot\overrightarrow{PQ}$ 的最小值为 \blankline
    \answer{$\frac{3}{4}$}
  \item 在正方体中随机取三条棱，它们两两异面的概率为 \blankline
    \answer{$\frac{2}{55}$}
  \item 在平面直角坐标系 $xOy$ 中，点集 $K = \{(x,y) \mid (|x|+|3y|-6)(|3x|+|y|-6) \leq 0 \}$ 所对应的平面区域的面积为 \blankline
    \answer{$24$}
  \item 设 $\omega$ 为正实数，若存在 $a$, $b$ ($\pi\leq a<b\leq 2\pi$), 使得 $\sin\omega a+\sin\omega b = 2$，则实数 $\omega$ 的取值范围是 \blankline
    \answer{$\omega\in[\frac{9}{4},\frac{5}{2}]\cup [\frac{13}{4},+\infty)$}
  \item 对四位数 $\overline{abcd}$ ($1\leq a\leq 9$, $0\leq b,c,d\leq 9$), 若 $a>b$, $b<c$, $c>d$，则称 $\overline{abcd}$ 为 $P$ 类数，若 $a<b$, $b>c$, $c<d$，则称 $\overline{abcd}$ 为 $Q$ 类数，
  用 $N(P)$ 和 $N(Q)$ 分别表示 $P$ 类数和 $Q$ 类数的个数，则 $N(P)-N(Q)$ 的值为 \blankline
    \answer{$285$}
\end{enumerate}

\section{解答题：本大题共 3 个小题，满分 56 分. 解答应写出文字说明、证明过程或演算步骤.}

\begin{enumerate}[resume]
  \item （本题满分 16 分）若实数 $a$, $b$, $c$ 满足 $2^a+4^b = 2^c$,
    $4^a+2^b = 4^c$，求 $c$ 的最小值.
  \item （本题满分 20 分）设 $a_1$, $a_2$, $a_3$, $a_4$ 是四个有理数，使得 $\{ a_ia_j\mid 1\leq i<j\leq 4\} = \{24,-2,-\frac32,-\frac18,1,3\}$，求 $a_1+a_2+a_3+a_4$ 的值.
  \item （本题满分 20 分）在平面直角坐标系 $xOy$ 中，$F_1$, $F_2$ 分别是
    椭圆 $\frac{x^2}{2}+y^2 = 1$ 的左右焦点，设不经过焦点 $F_1$ 的直线 $l$
    与椭圆 $C$ 交于两个不同的点 $A$, $B$，焦点 $F_2$ 到直线 $l$ 的距离为 $d$.
    如果直线 $AF_1$, $l$, $BF_1$ 的斜率成等差数列，求 $d$ 的取值范围.
\end{enumerate}



\end{document}